\documentclass[11pt, a4paper]{article}

\usepackage[a4paper, top=2.5cm, bottom=2.5cm, left=2cm, right=2cm]{geometry}
\usepackage{amsmath}
\usepackage{amssymb}
\usepackage[colorlinks=true, linkcolor=blue, urlcolor=blue, citecolor=blue]{hyperref}

\title{Theory of Everything - Volume 21 \\ \Large The Arrow of Time}
\author{Omega Protocol}
\date{\today}

\begin{document}

\maketitle

\section*{Layman's Summary}
Why can you turn an egg into an omelet, but you can't turn an omelet back into an egg? Why does time only flow forward? Volume 21 proves that time is not a fundamental dimension like space; rather, the "flow" of time is just our macroscopic experience of quantum information settling into more stable configurations.

\section{The Thermodynamic Arrow (Axiom 24)}
At the microscopic level, the equations of quantum mechanics don't care whether time runs forward or backward. We establish that the macroscopic, one-way "flow" of time is identical to the gradient of relative entropy dissipation. "Time" is just the sensation of information scrambling. If information isn't changing, time doesn't exist.

\section{Time Asymmetry (Theorem)}
How does a symmetric microscopic world create a one-way macroscopic world? We prove that macroscopic irreversibility naturally emerges from the symmetric microscopic "modular flow." Because macroscopic systems can only track coarse-grained, fuzzy information, the underlying exact symmetries are hidden, creating the illusion of a one-way street.

\section{The Past Hypothesis (Theorem)}
For entropy to always increase, it must have started very low. The "Past Hypothesis" is the assumption that the universe began in a highly ordered state. We prove that this is not an arbitrary guess; a low-entropy initial state is a necessary topological boundary condition of the underlying mathematical algebra. The universe had to start ordered.

\end{document}
